% ID: SEL-D0058-OA-E03
% ESTADO: cotejada. Segunda pasada visual y borrador compilado revisados.
% Fuente: Selectividad/sel_2011_matematicas_aplicadas/221A-EXAMEN 5.pdf
% SHA-256: 1780d560801fb3001e146f12b3a37b29b566ca6576dc9cdafda07902d7dee982
% Página: 1; opción A; ejercicio original 3.

En un sistema de alarma, la probabilidad de que haya un incidente es $0.1$. Si éste se produce, la probabilidad de que la alarma suene es $0.95$. La probabilidad de que suene la alarma sin que haya incidente es de $0.03$.
\begin{enumerate}[label=\alph*)]
\item \textbf{(1.5 puntos)} \textquestiondown{}Cuál es la probabilidad de que suene la alarma?
\item \textbf{(1 punto)} Si ha sonado la alarma, calcule la probabilidad de que no haya habido incidente.
\end{enumerate}
